\section{Further Evaluation}
\label{app:evaluation-agenda}
The next experiments should resolve specific gaps exposed by the pilots.
First, declare the main complete-host contrasts and repeat them on separate
days with additional fresh processes, recording process calibration and
matched continuous/paced behavior at identical boundaries. Use the pilots to
choose observation duration before collection; keep future confirmation
separate from these descriptive results. A second architecture and independent
build would test portability rather than merely extending this host's trace.

Second, compare paced native replacements at the hardest complete-module
settings and choose a background sweep that actually brackets saturation.
Record wake lateness, block compute, process CPU and device underruns separately.
The distinction between computation and host/device reliability also appears
in audio benchmark methodology~\cite{skare2024,balasubramaniam2026}. Physical
audio interfaces and concurrent rendering are required for user-facing
reliability and repaint claims.

Third, test native sub-FFT leaves and cost-informed weights under matched
input/output semantics. Transform subdivision~\cite{wefers2015} and alternative
traversals~\cite{becoulet2021} are credible candidates, not measured improvements.
Time-distributed inverse synthesis and long-filter chains need their own
bounded buffering, normalization and delivery paths.

Worker/UI-thread analysis should retain capture, queueing, overflow and
consumption costs. Overlap-reusing methods require equivalent window and
accuracy targets~\cite{garrido2016,park2014,rafii2018,eleftheriadis2023}; their
stability must be checked over long streams~\cite{vanderbyl2016}. Finally,
concurrent displays should measure actual repaint and visible age in addition
to the existing headless snapshot observations. These comparisons extend the
execution contract while testing alternatives outside the present study.
